\providecommand{\tightlist}{\setlength{\itemsep}{0pt}\setlength{\parskip}{0pt}}
\section{DEG\_L / integral ideal content of all CM transports}\label{deg_l-integral-ideal-content-of-all-cm-transports}

Run DEG\_L\_20261121; author /root/lvalue\_bounds; GPT-6 Astra Pro; effort not exposed. This tests whether integral combinations of many reference isogenies can remove the index appearing in B19. The answer below concerns pushforwards with their actual multiplicities. It does not apply to replacing each pushforward by its reduced image.

\subsection{B23. Relative determinant content survives every integral combination}\label{b23.-relative-determinant-content-survives-every-integral-combination}

Let \(K\) be an imaginary quadratic field, let \(E\) be a CM field containing \(K\), and put \[
 [E:K]=g=2n.
\] Fix a CM type of \(E\) whose restriction to \(K\) has signature \((n,n)\). Suppose that the rational homology has rank one over \(E\). Let \[
 A_0=\mathbf C^g/\Phi(I_0),\qquad
 A=\mathbf C^g/\Phi(\Lambda),
\] where \(I_0\) is an invertible fractional \(\mathcal O_E\)-ideal and \(\Lambda\) is a full lattice stable under \(\mathcal O_K\). The target is not assumed to have integral \(\mathcal O_E\) action. Simplicity and primitivity of the CM type are unnecessary for the lattice assertions below.

Let \(W_K(A)\subset H^g(A,\mathbf Q)\) be the usual two-dimensional rational Weil summand, the sum of the two extreme \(K\)-eigenspaces after extending scalars. Define its actual integral intersection lattice by \[
 \mathscr W_A=W_K(A)\cap H^g(A,\mathbf Z),
 \qquad
 \mathscr W_0=W_K(A_0)\cap H^g(A_0,\mathbf Z).
\tag{B23.1}
\] On this rational Weil plane use the determinant \(K\)-module structure: \(\mathrm{GL}_K(H^1)\) acts by its determinant. Thus the geometric endomorphism \([a]^*\) acts by the determinant scalar \(a^g\), not by \(a\). We identify the rational cohomology spaces using the trace pairing \(\operatorname{Tr}_{E/\mathbf Q}(xy)\). This fixes the conjugation convention in the formulas below.

Put \[
 \begin{aligned}
 \Lambda^-&=\{x\in\Lambda:\mathcal O_E x\subset\Lambda\},
 &J&=[\Lambda:\Lambda^-],\\
 I&=(\Lambda:I_0)=\operatorname{Hom}_E(A_0,A),
 &\mathfrak d&=\frac{\det_{\mathcal O_K}\Lambda}
                       {\det_{\mathcal O_K}I_0},\\
 \mathfrak j&=\frac{\det_{\mathcal O_K}\Lambda^-}
                       {\det_{\mathcal O_K}\Lambda}.
 \end{aligned}
\tag{B23.2}
\] Determinant modules are rank-one projective modules; their ratios are fractional ideals of \(K\). In particular, \(\mathfrak j\) is integral and \[
 N_{K/\mathbf Q}(\mathfrak j)=J.
\]

The following formulas are exact, including at primes ramified in \(K\): \[
 \begin{aligned}
 \mathscr W_A&=\mathfrak d^{-1}\mathscr W_0,\\
 N_{E/K}(I)&=\mathfrak j\mathfrak d,\\
 \sum_{0\ne h\in\operatorname{Hom}_E(A_0,A)}h_*\mathscr W_0
 &=\overline{\mathfrak j}\,\mathscr W_A.
 \end{aligned}
\tag{B23.3}
\] Here the sum means the additive subgroup generated by all the indicated integral images; it is already an \(\mathcal O_K\)-submodule. Thus its index in the target Weil lattice is exactly \(J\).

For a fixed nonzero \(w_0\in\mathscr W_0\), set \[
 \Gamma_{w_0}=\sum_{h}\mathbf Z\,h_*w_0,
 \qquad
 c_0=[\mathscr W_0:\mathcal O_Kw_0].
\] Choose once and for all \(a\in\mathcal O_K\) such that \(a^g\notin\mathbf R\), and put \[
 C=[\mathcal O_K:\mathbf Z[a^g]].
\tag{B23.4}
\] Then \[
 \begin{aligned}
 \mathcal O_K\Gamma_{w_0}
 &=\overline{\mathfrak j}\,\mathfrak d^{-1}\mathcal O_Kw_0,\\
 C\,\mathcal O_K\Gamma_{w_0}
 &\subseteq\Gamma_{w_0}
 \subseteq\mathcal O_K\Gamma_{w_0},\\
 Jc_0&\le[\mathscr W_A:\Gamma_{w_0}]\le C^2Jc_0.
 \end{aligned}
\tag{B23.5}
\] The number \(C\) is fixed independently of the target lattice. Hence even the additive span of every transport of one fixed nonzero integral Weil class retains index of order \(J\).

Such an \(a\) exists: choose a nonreal \(\theta\in\mathcal O_K\) and take \(a=m+\theta\) with \(m\in\mathbf Z\). The imaginary part of \((m+\theta)^g\) is a nonzero real polynomial in \(m\), with leading coefficient \(g\operatorname{Im}\theta\), so only finitely many integers are excluded.

\subsubsection{Consequence for algebraic cycles and division integers}\label{consequence-for-algebraic-cycles-and-division-integers}

Assume the fixed integral generic Weil projector of B6 is available: \[
 P_0=Re_0\quad\text{on }H^g(A_0),\qquad
 P_A=Re_A\quad\text{on }H^g(A),
\] with the same positive integer \(R\). Let \(\xi\) be any integral linear combination of the cohomology classes of multiplicity-preserving pushes \(h_*Z\), where \(h:A_0\to A\) is integral and \(E\)-linear and \(Z\) is any integral codimension-\(n\) cycle on \(A_0\). The source cycles may vary without any degree bound. Then \[
 P_A\xi\in\overline{\mathfrak j}\,\mathscr W_A.
\tag{B23.6}
\] In particular, if \(t\in\mathscr W_A\) is nonzero and \(\xi=Nt\) for a positive integer \(N\), put \[
 c_t=[\mathscr W_A:\mathcal O_Kt].
\] There is the necessary divisibility and lower bound \[
 J\mid R^2N^2c_t,
 \qquad
 N\ge\frac{\sqrt J}{R\sqrt{c_t}}.
\tag{B23.7}
\] In a fixed generic Weil family with integral \(\mathcal O_K\) action, after fixing the integral local system and a nonzero flat integral section \(t\), the integer \(c_t\) is constant. Thus no integral combination of these reference transports can represent a fixed such class using uniformly bounded division integers as \(J\to\infty\).

This is an obstruction for the indicated construction. It is not an intrinsic lower bound over all algebraic cycles on \(A\).

\subsection{Proof}\label{proof}

\subsubsection{1. The actual integral Weil lattice is a determinant-lattice functor}\label{the-actual-integral-weil-lattice-is-a-determinant-lattice-functor}

Let \(M\) be a full \(\mathcal O_K\)-lattice of rank \(g\) in a \(K\)-vector space \(V\). Let \(W(V)\) be its two-dimensional rational extreme exterior summand in \(\bigwedge^g_{\mathbf Q}V\). Define \[
 L(M)=W(V)\cap\bigwedge^g_{\mathbf Z}M.
\] Fix a rational prime \(p\). The ring \(\mathcal O_K\otimes\mathbf Z_p\) is a product of discrete valuation rings, and \(M\otimes\mathbf Z_p\) is free of rank \(g\) over this ring. Under a chosen local basis the two-dimensional rational summand is the determinant representation of \(\mathrm{GL}_g(K\otimes\mathbf Q_p)\).

The integral intersection in the standard local lattice is stable under the full ring \(\mathcal O_K\otimes\mathbf Z_p\). Indeed, for any \(b\) in that ring the \(\mathcal O_K\otimes\mathbf Z_p\)-linear endomorphism with matrix \(\operatorname{diag}(b,1,\ldots,1)\) preserves the standard lattice. Its ordinary \(g\)th exterior power preserves the integral exterior lattice and acts on the extreme summand by the determinant scalar \(b\). This argument works for nonunits and for the split product ring as well. The standard intersection is therefore a rank-one projective module over the maximal order.

If \(T\) is a \(K\otimes\mathbf Q_p\)-linear change of lattice, its ordinary exterior power takes the full integral exterior lattice isomorphically onto that of the new lattice and preserves the rational extreme summand. Consequently \[
 L(TM_p)=\det_K(T)L(M_p).
\tag{B23.8}
\] The possible scalar ideal distinguishing the standard intersection from the abstract determinant module is the same before and after the change of lattice, so it cancels in (B23.8). No unramifiedness assertion or division by the discriminant of \(K\) is being used.

These local identities glue to the corresponding determinant-ideal transformation law over \(\mathcal O_K\). For \(H^1\) one takes the integral \(\mathbf Z\)-dual of homology. Under the fixed trace pairing, taking this dual reverses the relative determinant ideal: locally a homology change matrix \(T\) changes its dual lattice by \((T^{-1})^{\mathsf t}\). The trace-dual factor of \(\mathcal O_K\) is common to source and target and again cancels. Since integral cohomology of an abelian variety is the exterior algebra on \(H^1(A,\mathbf Z)\), we obtain \[
 \mathscr W_A=\mathfrak d^{-1}\mathscr W_0.
\] This also proves that these actual integral intersection lattices are stable under the determinant \(\mathcal O_K\) action, a fact stronger than stability merely under the geometric scalars \(a^g\).

\subsubsection{2. Relative ideal norms and the action of a pushforward}\label{relative-ideal-norms-and-the-action-of-a-pushforward}

The B19 Hom ideal identity gives \(II_0=\Lambda^-\). Taking determinant modules over \(\mathcal O_K\) yields \[
 N_{E/K}(I)=\frac{\det_{\mathcal O_K}(II_0)}
                    {\det_{\mathcal O_K}I_0}
 =\mathfrak j\mathfrak d.
\tag{B23.9}
\] Here the determinant norm formula can be checked prime by prime: an \(\mathcal O_E\) prime \(\mathfrak P\) above \(\mathfrak p\) changes the determinant valuation over \(\mathcal O_K\) by its residue degree \(f(\mathfrak P/\mathfrak p)\). This is exactly the standard definition of the relative ideal norm. Similarly, local elementary divisors show that the numerical norm of \(\mathfrak j\) is the additive index \([\Lambda:\Lambda^-]=J\).

Let \[
 \nu=N_{K/\mathbf Q}(\mathfrak d)>0.
\] It is the ratio of the real covolumes of \(\Lambda\) and \(I_0\). For \(0\ne\alpha\in I\), let \(h_\alpha\) be the corresponding isogeny and write \(s_\alpha=N_{E/K}(\alpha)\). On rational Weil cohomology, \[
 h_\alpha^*=s_\alpha,\qquad
 \deg h_\alpha=\nu^{-1}s_\alpha\overline{s_\alpha}.
\] The identity \(h_{\alpha*}h_\alpha^*=(\deg h_\alpha)\mathrm{id}\) therefore gives \[
 h_{\alpha*}=\nu^{-1}\overline{s_\alpha}.
\tag{B23.10}
\] These are middle-degree cohomological pushforwards with their actual multiplicities.

The \(\mathcal O_K\)-ideal generated by \(N_{E/K}(\alpha)\) for \(\alpha\in I\) is exactly \(N_{E/K}(I)\). One inclusion follows from \((\alpha)\subseteq I\). For the reverse inclusion, fix a prime \(\mathfrak p\) of \(K\). By the Chinese remainder theorem for the invertible ideal \(I\), choose an element \(\alpha\in I\) attaining its minimal ideal valuation at every prime \(\mathfrak P\) of \(E\) above \(\mathfrak p\). Its norm attains precisely the valuation of \(N_{E/K}(I)\) at \(\mathfrak p\). This proves equality after localizing at every prime, hence equality of the ideals. Fractional ideals reduce to integral ideals by multiplication by a fixed nonzero denominator.

Each image \(h_*\mathscr W_0\) is an \(\mathcal O_K\)-module. Using (B23.9)--(B23.10), their additive sum is \[
 \begin{aligned}
 \sum_hh_*\mathscr W_0
 &=\nu^{-1}\overline{N_{E/K}(I)}\,\mathscr W_0\\
 &=\nu^{-1}\overline{\mathfrak j}\,\overline{\mathfrak d}\,\mathscr W_0\\
 &=\overline{\mathfrak j}\,\mathfrak d^{-1}\mathscr W_0
 =\overline{\mathfrak j}\,\mathscr W_A.
 \end{aligned}
\tag{B23.11}
\] The third equality uses \(\mathfrak d\overline{\mathfrak d}=(\nu)\) in the imaginary quadratic field. It proves all of (B23.3).

There is also a geometric factorization explaining the common content. Every \(h\) factors through \[
 A^-=\mathbf C^g/\Phi(\Lambda^-)
 \xrightarrow{q}A,
 \qquad\deg q=J.
\] The first part of the proof gives \(\mathscr W_{A^-}=\mathfrak j^{-1}\mathscr W_A\). In the common rational cohomology coordinates, \(q^*\) is the identity and \(q_*=J\). Therefore \[
 q_*\mathscr W_{A^-}=J\mathfrak j^{-1}\mathscr W_A
 =\overline{\mathfrak j}\,\mathscr W_A.
\tag{B23.12}
\] Thus the common sublattice arises before any choice of short isogeny, norm estimate, or theta input.

\subsubsection{3. A fixed source vector and its additive span}\label{a-fixed-source-vector-and-its-additive-span}

The same calculation with the principal rank-one submodule \(\mathcal O_Kw_0\) gives the first formula in (B23.5), of index \(Jc_0\) in \(\mathscr W_A\).

For every \(\alpha\in I\) one also has \(a\alpha\in I\), and (B23.10) gives \[
 h_{a\alpha*}w_0=\overline a^{\,g}h_{\alpha*}w_0.
\] Consequently \(\Gamma_{w_0}\) is a module over \(\mathbf Z[\overline a^{\,g}]\). The latter order has index \(C\) in \(\mathcal O_K\), so \[
 C\mathcal O_K\subseteq\mathbf Z[\overline a^{\,g}].
\] It follows that \(C\mathcal O_K\Gamma_{w_0}\subseteq\Gamma_{w_0}\). The other inclusion is tautological. Since all these lattices have rank two over \(\mathbf Z\), multiplication by \(C\) has index \(C^2\). This proves the two-sided index bound in (B23.5).

\subsubsection{4. Projection of integral cycles and the division bound}\label{projection-of-integral-cycles-and-the-division-bound}

The homomorphisms \(h\) are \(K\)-linear. They commute on cohomology with the fixed polynomial projector operators, including pushforwards: \[
 P_Ah_*=h_*P_0.
\] For each integral source cycle \(Z\), the class \(P_0[Z]\) belongs to \(\mathscr W_0\). Formula (B23.3) therefore proves (B23.6) for every integral linear combination, regardless of the degrees or number of source cycles used.

If \(\xi=Nt\) lies in the Weil plane, then \(P_A\xi=RNt\). Its membership in \(\overline{\mathfrak j}\mathscr W_A\) implies the lattice inclusion \[
 \mathcal O_K(RNt)\subseteq\overline{\mathfrak j}\mathscr W_A.
\] Taking additive indices in the rank-two lattice \(\mathscr W_A\) gives \[
 [\mathscr W_A:\mathcal O_K(RNt)]
 =R^2N^2c_t
 =J[\overline{\mathfrak j}\mathscr W_A:\mathcal O_K(RNt)].
\] This proves both assertions of (B23.7). The determinant \(K\) action, integral Weil intersection lattice, and a fixed integral section \(t\) are flat in the indicated generic Weil family. Their finite index \(c_t\) is therefore locally constant and hence constant on a connected base after fixing the local system. QED.

\subsection{More precise ideal divisibility}\label{more-precise-ideal-divisibility}

The necessary projected-class condition is more precise than its numerical norm. Define \[
 \mathfrak b_t=\{b\in\mathcal O_K:bt\in\overline{\mathfrak j}\mathscr W_A\},
 \qquad
 \mathfrak b_t\cap\mathbf Z=m_t\mathbf Z,
\tag{B23.13}
\] where \(m_t\) is positive. Then \(RNt\in\overline{\mathfrak j}\mathscr W_A\) is equivalent to \[
 \frac{m_t}{\gcd(m_t,R)}\mid N.
\tag{B23.14}
\] This is an exact test for membership in the common projected lattice. It is not a sufficiency assertion for algebraicity, the full cohomology class, or a prescribed effective cycle.

For example, if \(t\) generates \(\mathscr W_A\) over \(\mathcal O_K\), then \(\mathfrak b_t=\overline{\mathfrak j}\). If \(\mathfrak j=(m)\), the projected membership condition requires \(m/\gcd(m,R)\mid N\), and \(J=m^2\). If instead \(\mathfrak j=\mathfrak p^r\) for one of the primes above a split rational prime \(p\), then \(J=p^r\) and \(m_t=p^r\). For \(p\nmid R\), the condition requires \(p^r\mid N\). These are conditional computations for the displayed determinant ideals, not claims that every such ideal occurs in a prescribed polarized moduli family.

\subsection{A fixed nonmaximal quadratic order}\label{a-fixed-nonmaximal-quadratic-order}

The exact formulas above assume integral \(\mathcal O_K\) action. If the fixed generic action is only by an order \(\mathcal O\subset\mathcal O_K\), choose a fixed positive integer \(f\) with \(f\mathcal O_K\subset\mathcal O\). Set \[
 \Lambda^+=\mathcal O_K\Lambda,
 \qquad A^+=\mathbf C^g/\Phi(\Lambda^+),
 \qquad u:A\to A^+.
\] Then \(f\Lambda^+\subset\Lambda\), so \(q=\deg u=[\Lambda^+:\Lambda]\le f^{2g}\). Let \((\Lambda^+)^-\) denote its maximal \(\mathcal O_E\) core and put \(J^+=[\Lambda^+:(\Lambda^+)^-]\). There are inclusions \[
 f(\Lambda^+)^-\subseteq\Lambda^-\subseteq(\Lambda^+)^-.
\] If \(s=[(\Lambda^+)^-:\Lambda^-]\), then \(s\le f^{2g}\) and \[
 qJ=sJ^+,
 \qquad J^+\ge J/f^{2g}.
\tag{B23.15}
\]

Apply the proved maximal-order theorem to every composite \(uh\). If an integral combination of the original pushes has class \(Nt\), its image under \(u_*\) has class \(N t^+\), where \(t^+=u_*t\ne0\) is integral. On a family with fixed integral \(\mathcal O\) action and fixed flat \(t\), the saturation by \(\mathcal O_K\), the finite isogeny \(u\), the induced Weil lattices, and \(t^+\) are all determined by the same fixed rational local system. In particular \[
 c_{t^+}=[\mathscr W_{A^+}:\mathcal O_Kt^+]
\] is fixed. The projector commutes with \(u\) and keeps denominator \(R\). Thus \[
 N\ge\frac{\sqrt{J^+}}{R\sqrt{c_{t^+}}}
 \ge\frac{\sqrt J}{R f^g\sqrt{c_{t^+}}}.
\tag{B23.16}
\] This supplies a uniform comparison for a fixed quadratic order. It does not silently replace a varying quadratic order by its maximal order at no cost.

\subsection{Reduced images from an arbitrary bounded source family}\label{reduced-images-from-an-arbitrary-bounded-source-family}

Assume now that \(A_0\) is simple, as in B21--B22. Fix positive integers \(D_0\) and \(k\), with \(L_0^{\otimes k}\) very ample, and put \[
 \begin{aligned}
 S&=k^{n(n+1)}D_0^{n+1}\in\mathbf Z_{>0},\\
 L_S&=\operatorname{lcm}(1,2,\ldots,S).
 \end{aligned}
\] The sharper intrinsic result B22a permits the smaller choice \[
 S=\max\left\{1,\left\lfloor
 \frac{D_0^2}{4(n!)^2\chi(L_0)}\right\rfloor\right\}.
\] Every assertion below holds with either choice, or with any other positive integer upper bound for the stabilizer orders in this source family. Let \(\xi\) be any integral linear combination of reduced-image classes \[
 [h(Z)_{\mathrm{red}}],
 \qquad
 Z\subset A_0\text{ integral of dimension }n,
 \qquad \deg_{L_0}Z\le D_0,
\] with arbitrary integral \(E\)-linear isogenies \(h:A_0\to A\). There is no bound on the number of these summands or their integer coefficients, and the source subvarieties may vary throughout the entire bounded-degree family. If \(\xi=Nt\) for a nonzero integral target Weil class \(t\), then in the maximal-order setting \[
 J\mid R^2L_S^2N^2c_t.
\tag{B23.17}
\] Consequently the two order conventions give respectively \[
 \begin{aligned}
 N&\ge\frac{\sqrt J}{R L_S\sqrt{c_t}}
 &&\text{for integral }\mathcal O_K\text{ action},\\
 N&\ge\frac{\sqrt J}{R L_S f^g\sqrt{c_{t^+}}}
 &&\text{for the fixed order }\mathcal O.
 \end{aligned}
\tag{B23.18}
\]

To prove this, B20 and B22 give, for each summand, \[
 h_*Z=\delta_hh(Z)_{\mathrm{red}},\qquad
 \delta_h=|\ker h\cap\operatorname{Stab}(Z)|\le S.
\] Thus \(\delta_h\) divides \(L_S\). Multiplying the entire proposed relation by \(L_S\) expresses \(L_S\xi=L_SN t\) as an integral linear combination of multiplicity-preserving source pushes. Applying (B23.6)--(B23.7) proves (B23.17) and the first bound in (B23.18). For the fixed nonmaximal order, apply the same argument after pushing this multiplied relation through \(u:A\to A^+\). It gives \(J^+\mid R^2L_S^2N^2c_{t^+}\), and (B23.15) proves the second bound. No new multiplicity denominator is introduced by this final pushforward.

The sharper condition is \(m_t\mid RL_SN\) in the notation (B23.13). In particular, for a rational prime \(p>S\) one has \(v_p(L_S)=0\), so division by generic mapping multiplicities removes none of the \(p\)-primary condition \(v_p(N)\ge\max(0,v_p(m_t)-v_p(R))\). Integral greatest-common-divisor extraction across arbitrarily many reduced images therefore still cannot produce a fixed flat target class with bounded \(N\) along unbounded \(J\), provided the source degrees remain bounded. All constants refer to the fixed source polarization, quadratic order, integral local system and class. This assertion still imposes no intrinsic lower bound on cycles obtained by a different construction.

\subsection{Scope and the remaining route}\label{scope-and-the-remaining-route}

Taking greatest common divisors across all \(E\)-linear isogenies removes individual principal relative norm factors only down to the relative norm ideal of the full Hom lattice. Its remaining target content is the conjugate determinant ideal \(\overline{\mathfrak j}\), of norm \(J\). The same obstruction applies to integral sums described as traces if every summand is one of the indicated multiplicity-preserving reference transports. Operations leaving this class of summands require a separate argument.

Replacing \(h_*Z=\delta_h h(Z)_{\mathrm{red}}\) by the reduced image divides its class by \(\delta_h\). Such a division can leave the sublattice in (B23.6); this is precisely why B20--B22 separately control stabilizers and generic multiplicities. The present result imposes no new obstruction on arbitrary target cycles or new deformation constructions. It is also distinct from a single full-CM tensor coefficient in the product construction of 032: applying it there requires identifying an appropriate generic rank-two \(K\)-Weil summand and its integral lattice.

\subsubsection{Exact classical inputs and reading scope}\label{exact-classical-inputs-and-reading-scope}

\begin{itemize}
\tightlist
\item
  J.S. Milne, \emph{Algebraic Number Theory}, version 3.08, 19 July 2020, section 4, printed pp.68--70, definition of the relative ideal norm and Propositions 4.1--4.2. Official source \url{https://www.jmilne.org/math/CourseNotes/ANT.pdf}. The definition, compatibility with principal norms, ideal-index formula, and the complete adjoining short proofs were read. These ideal-theoretic facts are imported; the determinant-lattice and transport applications above are proved here.
\item
  Integral cohomology of a complex abelian variety as the exterior algebra on its integral first cohomology, and finite-isogeny pull-push, are the same inputs already used and checked in B7, B18 and B20. The local maximal-order lattice argument above explicitly handles ramified and split primes; no unrecorded saturation denominator is assumed.
\item
  B6 supplies the common integral projector; B19 supplies \(II_0=\Lambda^-\) and the identification of integral \(E\)-linear Hom with the ideal \(I\). Simplicity is not used in these algebraic identities, as also checked in T6.17.
\end{itemize}

This proof has been read in full, including the determinant action versus geometric scalar distinction, the actual integral intersection lattice, the relative norm ideal generated by all element norms, and the fixed-order comparison. No proof-assistant verification is claimed.
